% ==============================================================================
% Chapitre 1 : Introduction
% Contexte, objectifs, périmètre, méthodologie et organisation du document
% ==============================================================================

\chapter{Introduction}

\section{Contexte du projet}

\subsection{Cadre académique}

Ce dossier d'architecture s'inscrit dans le cadre du projet annuel de
2\textsuperscript{e} année à l'\textbf{ESGI} (École Supérieure de Génie
Informatique) de Toulouse. Le projet fil rouge se déroule sur l'année
2025--2026 et constitue la principale évaluation pratique du cursus. Il
réunit une équipe de quatre étudiants dont la mission est de réaliser un
projet informatique complet, de la phase de conception jusqu'au
déploiement, en passant par le développement, les tests et la mise en
production.

L'objectif pédagogique est de confronter les étudiants à l'ensemble des
étapes du cycle de vie d'un logiciel, dans des conditions proches de
celles rencontrées en entreprise. Le projet doit non seulement
fonctionner techniquement, mais aussi respecter des contraintes de
qualité, de sécurité et de maintenabilité.

\subsection{Entreprise fictive : LockBits}

\textbf{LockBits} est une entreprise fictive de cybersécurité créée dans
le cadre de ce projet. Elle propose une solution \textbf{EDR} (Endpoint
Detection \& Response) open-source, associée à un portail client pour la
gestion des clients.

La solution EDR permet aux entreprises de détecter les menaces sur leurs
postes de travail (endpoints) et d'y répondre de manière orchestrée.
Elle se compose de deux éléments principaux~: d'une part un serveur
central qui collecte et analyse les données de sécurité, d'autre part un
agent déployé sur chaque machine à protéger. Le portail client est un composant essentiel de la solution, permettant
aux clients de gérer leur sécurité en temps réel.

LockBits se positionne comme une alternative open-source aux solutions
EDR propriétaires du marché. Le choix de l'open-source permet à la fois
de démontrer la transparence de la solution et de faciliter son adoption
par les entreprises soucieuses de maîtriser leur outillage de sécurité.

\subsection{Enjeux et défis}

Le projet LockBits couvre l'ensemble du cycle de vie logiciel, ce qui
introduit plusieurs enjeux majeurs.

Sur le plan technique, l'architecture est complète et hétérogène~: un
backend en Python/FastAPI, un frontend web en PHP, une base de données
relationnelle, un système de monitoring, et une chaîne d'intégration et
de déploiement continus. Chaque composant doit s'intégrer harmonieusement
avec les autres.

Sur le plan de la sécurité, le produit manipule des données système
sensibles~: processus en cours d'exécution, connexions réseau, fichiers
suspects. La collecte et le traitement de ces informations imposent des
mesures de protection strictes, tant au niveau du transport que du
stockage. La solution doit respecter les principes de moindre
privilège et de défense en profondeur.

Sur le plan de la disponibilité, l'ensemble de l'application est
critique pour les clients. Le portail client et GLPI doivent
fonctionner pour que les clients reçoivent les alertes de sécurité.
Le système doit donc être résilient et permettre des mises à jour
sans interruption de service.

\section{Objectifs du projet}

\subsection{Objectifs techniques}

Le projet LockBits poursuit sept objectifs techniques principaux~:

\begin{enumerate}
  \item \textbf{Serveur EDR}~: développer un serveur en FastAPI/Python
    exposant une API REST pour la collecte des données des agents, le
    stockage des rapports et l'analyse de fichiers suspects via
    VirusTotal.
  \item \textbf{Agent EDR}~: créer un agent déployable sur les endpoints
    Linux, disponible sous forme de binaire portable et d'image Docker,
    capable de collecter les processus, les connexions réseau et les
    événements de modification de fichiers (FIM).
  \item \textbf{Site client web}~: réaliser un portail client en
    PHP/Apache avec authentification, tableau de bord, gestion des
    clients et système de tickets de support.
  \item \textbf{Chatbot intelligent}~: intégrer un assistant virtuel au
    site client, alimenté par l'API Claude d'Anthropic, pour répondre
    aux questions courantes des utilisateurs.
  \item \textbf{Pipeline CI/CD}~: automatiser l'intégration et le
    déploiement via GitHub Actions, avec des étapes de lint, test, build
    multi-architecture, scan de sécurité et déploiement.
  \item \textbf{Monitoring et observabilité}~: mettre en place une
    infrastructure de supervision avec Prometheus, Grafana et Loki pour
    assurer la visibilité sur l'état du système.
  \item \textbf{Déploiement automatisé}~: fournir un script d'installation
    one-liner (install.sh) pour déployer l'intégralité de la stack en
    quelques commandes.
\end{enumerate}

\subsection{Objectifs pédagogiques}

Au-delà des objectifs techniques, le projet vise des objectifs
pédagogiques précis. Les étudiants doivent mettre en pratique les
compétences acquises durant leur formation~: développement logiciel,
conception d'architecture, administration système, sécurité
informatique et gestion de bases de données.

La gestion de projet occupe une place centrale. L'équipe doit organiser
son travail, répartir les tâches et tenir des réunions hebdomadaires.
Cette organisation reproduit les méthodes agiles utilisées en
entreprise, sans suivre de sprints formels.

Le projet introduit également les outils \textbf{DevOps} modernes~:
Docker pour la conteneurisation, GitHub Actions pour l'intégration
continue, Prometheus et Grafana pour le monitoring. Ces compétences sont
aujourd'hui essentielles dans le monde professionnel.

\section{Périmètre et contraintes}

\subsection{Périmètre fonctionnel}

Le périmètre fonctionnel du projet LockBits couvre les domaines
suivants~:

\begin{itemize}
  \item \textbf{Détection et réponse sur endpoints}~: collecte en
    continu des processus actifs, des connexions réseau et des
    modifications du système de fichiers. L'agent peut fonctionner en
    mode scan (analyse ponctuelle) ou en mode monitor (surveillance en
    temps réel).
  \item \textbf{Portail client}~: inscription et connexion des clients,
    tableau de bord avec les alertes de sécurité, gestion du profil et
    abonnement, système de tickets de support.
  \item \textbf{Chatbot assistant}~: assistant virtuel intégré au portail
    pour répondre aux questions fréquentes et orienter les utilisateurs
    vers les ressources appropriées.
  \item \textbf{Mise à jour automatique}~: les composants peuvent se
    mettre à jour automatiquement via un mécanisme de cron, garantissant
    que les clients bénéficient toujours des dernières correctifs de
    sécurité.
  \item \textbf{Intégration VirusTotal}~: analyse des fichiers suspects
    remontés par les agents via l'API publique de VirusTotal, avec
    enrichissement automatique des rapports.
  \item \textbf{Intégration GLPI}~: synchronisation avec GLPI pour la
    gestion des tickets IT, avec possibilité de SSO (Single Sign-On)
    entre le portail client et l'outil de ticketing.
\end{itemize}

\subsection{Contraintes techniques}

La stack technique est imposée par le cahier des charges du projet~:

\begin{itemize}
  \item \textbf{EDR Server}~: Python avec FastAPI, base de données
    SQLite pour le stockage local des rapports.
  \item \textbf{Site client}~: PHP 8+ avec Apache, base de données MySQL
    8.0.
  \item \textbf{Conteneurisation}~: Docker avec Docker Compose pour
    l'orchestration des services, chaque composant possédant son propre
    \texttt{Dockerfile}.
  \item \textbf{CI/CD}~: GitHub Actions pour l'automatisation des
    pipelines, GitHub Container Registry (GHCR) pour le stockage des
    images.
  \item \textbf{Déploiement}~: Docker Compose en production, avec
    Dokploy comme plateforme d'hébergement. La branche \texttt{develop}
    est déployée en pré-production et la branche \texttt{main} en
    production.
  \item \textbf{Architecture matérielle}~: support des architectures
    \texttt{linux/amd64} et \texttt{linux/arm64}.
\end{itemize}

\subsection{Contraintes organisationnelles}

L'équipe est composée de quatre étudiants, ce qui impose une répartition
claire des responsabilités. Le projet se déroule sur environ huit mois,
de septembre à juillet. Les composants sont développés en parallèle via
des sous-modules Git, chaque sous-module étant un dépôt indépendant avec
son propre cycle de développement.

Le travail est organisé avec des réunions hebdomadaires de suivi, sans
sprints formels. Cette organisation permet de détecter rapidement les
blocages et d'ajuster les priorités.

\section{Méthodologie et approche}

\subsection{Approche itérative}

Le développement de LockBits suit une approche itérative et
incrémentale. Chaque composant est développé indépendamment dans son
propre dépôt, avec des cycles de validation courts. Une fois qu'un
composant atteint un état fonctionnel, il est intégré à la stack
globale via le dépôt principal qui orchestre l'ensemble via les
sous-modules Git.

Cette approche présente plusieurs avantages. Elle permet de paralléliser
le travail entre les membres de l'équipe. Elle réduit les risques
d'intégration en détectant les incompatibilités tôt. Elle offre une
visibilité continue sur l'avancement, chaque sous-module produisant des
livrables intermédiaires.

Chaque itération comprend les phases suivantes~:
définition des objectifs, développement, tests, intégration et revue.
Les retours d'une itération alimentent la planification de la suivante.

\subsection{Principes de conception}

La conception de l'architecture repose sur cinq principes fondamentaux~:

\begin{enumerate}
  \item \textbf{Modularité}~: chaque composant est indépendant et
    communique via des interfaces clairement définies (API REST). Les
    composants sont développés dans des dépôts Git distincts, liés par
    des sous-modules.
  \item \textbf{Conteneurisation}~: tous les services sont empaquetés
    dans des images Docker multi-architecture, garantissant la
    reproductibilité des déploiements et la portabilité entre les
    environnements.
  \item \textbf{Infrastructure as Code}~: l'ensemble de la chaîne
    d'intégration et de déploiement est automatisée via GitHub Actions.
    La configuration des services est versionnée et déployée de manière
    reproductible.
  \item \textbf{Sécurité dès la conception}~: les principes de sécurité
    sont intégrés à chaque étape du pipeline~: scan des vulnérabilités
    avec Trivy, gestion des secrets via variables d'environnement,
    rate limiting sur les endpoints exposés, validation des entrées.
  \item \textbf{Observabilité}~: le système produit des métriques, des
    logs et des alertes permettant de diagnostiquer les problèmes et
    d'anticiper les dégradations. Prometheus collecte les métriques,
    Grafana les visualise et Loki centralise les logs.
\end{enumerate}

\section{Organisation du document}

Ce dossier d'architecture est organisé en sept chapitres~:

\begin{itemize}
  \item \textbf{Chapitre 2~: Architecture globale}~: présente la vue
    d'ensemble du système, les diagrammes d'architecture, les choix
    techniques fondamentaux et les flux de données entre les composants.
  \item \textbf{Chapitre 3~: Composants applicatifs}~: détaille chaque
    composant du système~: le serveur EDR, l'agent EDR, le portail
    client, l'intégration GLPI et le chatbot.
  \item \textbf{Chapitre 4~: Infrastructure et déploiement}~: décrit
    l'infrastructure Docker, le pipeline CI/CD, le monitoring et les
    mesures de sécurité.
  \item \textbf{Chapitre 5~: Gouvernance}~: présente l'organisation de
    l'équipe, la méthodologie de travail, la gestion des risques et les
    jalons du projet.
  \item \textbf{Chapitre 6~: Extensions et évolutions}~: explore les
    pistes d'amélioration et la roadmap du projet.
  \item \textbf{Chapitre 7~: Annexes}~: regroupe le glossaire, les
    références techniques et les commandes utiles.
\end{itemize}
